% Generated draft: requires mathematical and source review.
\documentclass[11pt]{article}
\usepackage[margin=25mm]{geometry}
\usepackage{fontspec}
\setmainfont{Latin Modern Roman}
\usepackage{amsmath,amssymb,mathtools,unicode-math}
\setmathfont{Latin Modern Math}
\usepackage{enumitem,needspace,xurl,hyperref}
\setlength{\parindent}{0pt}
\setlength{\parskip}{5pt}
\setlength{\emergencystretch}{3em}
\newcommand{\N}{\mathbb{N}}
\newcommand{\Z}{\mathbb{Z}}
\newcommand{\Q}{\mathbb{Q}}
\newcommand{\R}{\mathbb{R}}
\newcommand{\C}{\mathbb{C}}
\newcommand{\F}{\mathbb{F}}
\newcommand{\A}{\mathbb{A}}
\newcommand{\PP}{\mathbb{P}}
\newcommand{\E}{\mathbb{E}}
\newcommand{\eps}{\varepsilon}
\newcommand{\dd}{\,\mathrm{d}}
\begin{document}
\section*{UnsolvedMath ID: 2302072}
\textbf{Source problem code:} catalogue entry
\subsection{Definitions and mathematical statement}
For a homogeneous linear equation of order $n$, a fundamental system is a collection of $n$ linearly independent solutions, equivalently with nonzero Wronskian. Finitely many nonreal zeros means that the zeros in $\mathbb C\setminus\mathbb R$ form a finite set; there is no finiteness requirement on real zeros.

\textbf{Source target (S. Hellerstein and J. Rossi).} Replace finitely many zeros by finitely many nonreal zeros in the reported Frank characterization of polynomial-coefficient differential equations having a fundamental system of solutions with that property. Does the same equivalence with reduction to constant coefficients continue to hold? The source says the unresolved range may be taken as $n\geq3$, citing Hellerstein--Rossi and Gundersen for the lower-order background.

Two formula defects prevent a unique precise equivalence. The extracted operator is
\[
L_n(w)=w^{(n)}+a_{n-1}(z)w^{(n-1)}+\cdots+a_0(z),
\]
with polynomial $a_j$, but its final term lacks a factor $w$. A standard homogeneous linear version would instead have $\sum_{j=0}^{n-1}a_j(z)w^{(j)}$, where $w^{(0)}=w$. This is an explicitly identified possible repair, not the supplied equation as written. The extracted transformation is
\[
w(z)=\exp\big(q(z)u(z)\big),
\]
where $q$ is a polynomial. That is different from the usual linear gauge substitution $w=e^{q}u$ and does not by itself specify a linear conjugation of differential operators.

Once the intended operator and substitution are recovered, the constant-coefficient target would need to be specified as an equation such as $u^{(n)}+\sum_{j=0}^{n-1}c_ju^{(j)}=0$ with constants $c_j$. One must also fix whether the finite-nonreal-zero condition concerns the particular displayed fundamental system or existence of some such system. No repaired theorem or equivalence is asserted here.
\subsection{Short English statement}
Does Frank's finite-zero reduction theorem extend to solutions with only finitely many nonreal zeros? The extracted equation and transformation are malformed, and the fundamental-system quantifier also needs clarification.
\subsection{Sources}
\begin{itemize}
\item Hayman, \emph{Research Problems in Function Theory} (2018), Problem 2.72, attributed to S. Hellerstein and J. Rossi, \url{https://arxiv.org/abs/1809.07200}; source-TeX extraction, accessed 2026-07-29. The source label 2.72 is retained as the local equation/problem reference.
\item Frank, reported finite-zero characterization; Hellerstein and Rossi, and Gundersen, lower-order background. The extracted citation entries give no titles, years, or theorem locations for these results.
\item Local registry \texttt{data/unsolvedmath/problems.json}, permanent ID $2302072$, source code catalogue entry, \emph{Research Problems in Function Theory --- Problem 2.72}. The record labels this open and OPEN-TRIAGE, reporting no recent resolution located, with NEEDS\_REVIEW source-status and rights qualifications.
\end{itemize}
\end{document}
